\chapter{Numerical Methods}

\section{Introduction}
For the better part of the course we focused on classical analytic solutions to geophysical PDE's.  You may have noticed that the situations that we have imagined to solve those PDE's were very specific.  While these solutions taught us some fundamental insights into key geophysical concepts and processes, the solutions are of limited utility in many real world situations.  This limitation arises from the true 3-D (or even 4-D) variations in the Earth's properties due to the physical state and/or composition.  Unfortunately, for many of these cases, there may not be an analytic solution.  So how then can we solve these PDE's?

Another issue we have to confront in geophysics is being able to solve problems using observations to predict the properties of the Earth.  Everything we have done thus far assumes that we know the distribution of the Earth's properties, which we use to compute the geophysical response field.  Since the field is what we observe, we are trying to solve the inverse problem.

Both of these problems have a vast mathematical foundation and literature exploring their many facets, ideal choices, estimates of error and uncertainty, etc.  This chapter will explore a few of the numerical concepts commonly employed to address solutions to geophysical equations and inversion and is meant at but a taste of the utility of these methods.  For those of you that have now decided that geophysics is not interesting, or for you, these methods can be useful to nearly any aspect of geosciences.

